%% kapa: exam=1 round=35 year=2024 period=2 form=A subject=accounting subject_ko=회계학 seq=36 number=76
%% source: Q-Net 감정평가사 기출문제 (cst00309) artlSeq=5213229 file=raw/1008/5213229/2교시 문제지 원본.pdf
%% conversion: kordoc 4.24.51 -> tools/gen1.py (auto); review=proofread-sonnet; flags=tab
\begin{question}{76}
(주)감평은 결합공정을 거쳐 주산품 A, B와 부산품 F를 생산하여 주산품 A, B는 추가가공한 후 판매하고, 부산품 F의 회계처리는 생산시점에서 순실현가치법(생산기준법)을 적용한다. (주)감평의 당기 생산 및 판매 자료는 다음과 같다.
\par{\small\centering\begin{tabular}{|c|>{\centering\arraybackslash}p{0.2\linewidth}|>{\centering\arraybackslash}p{0.2\linewidth}|c|c|}\hline
구분 & 분리점 이후 추가가공원가 & 추가가공 후 단위당 판매가격 & 생산량 & 판매량 \tabularnewline\hline
A & ￦1,000 & ￦60 & 100단위 & 80단위 \tabularnewline\hline
B & 200 & 30 & 140 & 100 \tabularnewline\hline
F & 500 & 30 & 50 & 40 \tabularnewline\hline
\end{tabular}\par}
결합원가 ￦1,450을 분리점에서의 순실현가능가치 기준으로 각 제품에 배분할 때 주산품 A의 매출총이익은? (단, 기초 재고자산은 없다.)
\begin{choices}
  \item ￦2,714
  \item ￦2,800
  \item ￦2,857
  \item ￦3,714
  \item ￦3,800
\end{choices}
\end{question}
\answer{5}
